\DocumentMetadata{pdfstandard=A-2b,pdfversion=1.7,lang=pt-BR}
\documentclass{abntexto-ufc}

\ifdefined\UFCArticleCompatibilitySetup
  \ufcsetup{
    tipo = artigo,
    impressao = anverso,
    autor = {Maria Silva},
    titulo = {Perfil Normativo de Artigo Científico},
    subtitulo = {Fixture de Validação},
    variacao-titulo = {Normative Scientific Article Profile: Validation Fixture},
    data-submissao = {15 de agosto de 2026},
    data-aprovacao = {29 de agosto de 2026},
    curriculo-resumido-autor = {Pesquisadora na área de normalização acadêmica},
    vinculo-institucional-autor = {Universidade Federal do Ceará},
    contato-autor = {maria.silva@example.org},
    fonte = times,
    fonte-estrita = nao
  }
\else
  \ufcsetup{
    type = article,
    print-mode = single-sided,
    author = {Maria Silva},
    title = {Perfil Normativo de Artigo Científico},
    subtitle = {Fixture de Validação},
    title-variant = {Normative Scientific Article Profile: Validation Fixture},
    submission-date = {15 de agosto de 2026},
    approval-date = {29 de agosto de 2026},
    author-short-curriculum = {Pesquisadora na área de normalização acadêmica},
    author-affiliation = {Universidade Federal do Ceará},
    author-contact = {maria.silva@example.org},
    font = times,
    strict-font = false
  }
\fi

\ufcbibliografia{tests/fixtures/referencias-v2.bib}

\begin{document}

\ifdefined\UFCArticleCompatibilitySetup
  \imprimircabecalhoartigo
\else
  \ufcPrintArticleHeader
\fi

\imprimirresumo{tests/fixtures/article-summary.tex}
\imprimirabstract{tests/fixtures/article-abstract.tex}

\ifdefined\UFCArticleCompatibilitySetup
  \imprimirdatasartigo
\else
  \ufcPrintArticleDates
\fi

\section{Introdução}
O artigo científico exige uma organização contínua dos elementos, sem a quebra automática de página entre as seções primárias. Esta fixture verifica o comportamento básico do perfil implementado para a Universidade Federal do Ceará.

A introdução também confirma o recuo de primeira linha de dois centímetros e o espaçamento simples adotado pelo perfil, preservando o corpo tipográfico de doze pontos definido pela classe.

\section{Desenvolvimento}
O desenvolvimento utiliza o mesmo fluxo contínuo e mantém a numeração progressiva das seções. A implementação é deliberadamente isolada dos perfis de trabalhos acadêmicos e projetos de pesquisa já certificados.

A bibliografia continua sendo processada pela infraestrutura vigente da classe, com a edição corrente da norma de referências selecionada pela camada de compatibilidade existente.

\subsection{Validação do perfil}
O teste exercita as formas canônica em inglês e de compatibilidade em português para a seleção do perfil e para os metadados específicos do artigo.

\section{Considerações finais}
A fixture termina o corpo textual sem solicitar nova página antes das referências, permitindo verificar que o cabeçalho REFERÊNCIAS segue o fluxo contínuo previsto para artigos científicos.

\nocite{silva2020}
\imprimirreferencias

\end{document}
